% Sekcja końcowa — dołącz w main.tex: % Sekcja końcowa — dołącz w main.tex: % Sekcja końcowa — dołącz w main.tex: % Sekcja końcowa — dołącz w main.tex: \input{podsumowanie}
% Wymaga pakietów z main.tex: tcolorbox, enumitem

\section*{Podsumowanie i refleksja projektowa}

Celem projektu było zaprojektowanie i porównanie trzech wariantów asystenta uczelnianego AGH
opartego na generowaniu wspomaganym pobieraniem informacji (RAG): wyszukiwania wyłącznie wektorowego,
hybrydowego GraphRAG (graf Neo4j + wektory Chroma) oraz wieloetapowego pipeline'u agentowego (LangGraph).
Ocenę przeprowadziliśmy na zestawie 18 scenariuszy --- 13 pytań merytorycznych dotyczących rekrutacji
i studiów oraz 5 prób obejścia zabezpieczeń (manipulacja promptem, fałszywe źródła, żądania nielegalne).
Dla każdego pytania zarejestrowaliśmy odpowiedzi trzech systemów, czas reakcji oraz --- w wariancie agentowym ---
pełny ślad etapów przetwarzania (trace JSON).

\vspace{1em}

% ---------------------------------------------------------------------------
\subsection*{Wnioski z pracy zespołowej}

\begin{tcolorbox}[colback=blue!5!white,colframe=blue!75!black,title=Doświadczenia edukacyjne zespołu]
    \begin{itemize}[leftmargin=*, nosep]
        \item \textbf{Projektowanie bazy wiedzy hybrydowej.}
              Zrozumieliśmy, że połączenie danych strukturalnych (tabele CSV załadowane do Neo4j)
              z tekstem regulaminów (fragmenty PDF w Chroma) wymaga świadomego modelowania ontologii ---
              relacji między kierunkiem, progiem punktowym, przedmiotem maturalnym czy terminem rekrutacji.
              Jakość odpowiedzi zależy w równym stopniu od schematu grafu i jakości wektoryzacji dokumentów.

        \item \textbf{Brak uniwersalnie najlepszej architektury.}
              Wyniki benchmarku potwierdziły, że żaden wariant nie dominuje we wszystkich kategoriach zapytań.
              Wyszukiwanie wektorowe dobrze odtwarza fragmenty tekstu z zachowaną logiką słowną
              (np.\ operator „albo'' w tabeli przedmiotów), lecz zawodzi przy filtrowaniu numerycznym
              (progi poniżej 700 punktów) i zapytaniach wielowarunkowych (terminy konkretnego cyklu rekrutacji).
              GraphRAG precyzyjnie obsługuje zapytania relacyjne i arytmetyczne, natomiast płaski zapis encji
              w grafie może zniekształcić reguły logiczne źródłowe.

        \item \textbf{Orkiestracja wieloetapowa jako warstwa kontroli.}
              Budowa grafu stanów w LangGraph nauczyła nas rozdzielać odpowiedzialności między wyspecjalizowane
              moduły: strażnika bezpieczeństwa, analizator intencji, filtr kontekstu, generator odpowiedzi
              i audytora faktów. Każdy etap podejmuje decyzję warunkową (odmowa, doprecyzowanie, akceptacja),
              co odróżnia system agentowy od liniowego schematu \textit{pobierz\,--\,wygeneruj}.

        \item \textbf{Konieczność systemowej ewaluacji.}
              Opracowanie narzędzia porównawczego (\texttt{run\_compare.py}) i rejestrowanie śladów JSON
              uświadomiło nam, że ocena asystenta nie może ograniczać się do poprawności merytorycznej końcowej odpowiedzi.
              Istotne są także: etap, na którym system zawiódł, czas reakcji, zachowanie przy braku danych
              oraz sposób reagowania na próby manipulacji.

        \item \textbf{Praca zespołowa w projekcie inżynierskim.}
              Podzieliliśmy zakres na komplementarne obszary: przygotowanie i ingest danych, moduły pobierania
              informacji, zamrożone pipeline'y porównawcze, warstwę agentową oraz dokumentację z benchmarkiem.
              Integracja komponentów wymagała uzgodnienia wspólnych interfejsów stanu i formatów kontekstu,
              co stanowiło istotny element nauki inżynierii oprogramowania na rzeczywistym przypadku użycia.
    \end{itemize}
\end{tcolorbox}

\vspace{1em}

% ---------------------------------------------------------------------------
\subsection*{Zidentyfikowane słabości}

\begin{tcolorbox}[colback=red!5!white,colframe=red!50!black,title=Ograniczenia obecnej implementacji]
    \begin{itemize}[leftmargin=*, nosep]
        \item \textbf{Niejednorodna jakość odpowiedzi.}
              W zależności od kategorii pytania każdy wariant wykazuje inne słabości.
              Przykład: przy pytaniu o punkty za osiągnięcia sportowe wyszukiwanie wektorowe zwróciło fragment
              o punktach ECTS zamiast tabeli rekrutacyjnej; GraphRAG pomylił tabele (odznaki PTTK zamiast klas sportowych);
              system agentowy odrzucił błędny szkic w weryfikacji faktów, lecz użytkownik nie otrzymał pełnej,
              poprawnej odpowiedzi z listą dyscyplin.

        \item \textbf{Ograniczenia modelu grafowego.}
              Przy przedmiotach maturalnych na kierunku Automatyka i Robotyka graf zwrócił płaską listę encji
              bez operatora logicznego „albo'', co doprowadziło GraphRAG i RAG agentowy do błędnej interpretacji
              (wymaganie wszystkich przedmiotów naraz). Przy pytaniu o stypendium socjalne graf nie zwrócił
              żadnych wyników, a system oparł się wyłącznie na słabych fragmentach wektorowych.

        \item \textbf{Wysoki koszt obliczeniowy wariantu agentowego.}
              Pipeline agentowy wymaga od czterech do sześciu wywołań modelu językowego na jedno pytanie,
              co przekłada się na wielokrotnie wyższą latencję w porównaniu z wariantami liniowymi.
              W scenariuszach interaktywnych może to obniżać użyteczność systemu.

        \item \textbf{Zależność od jakości i aktualności danych wejściowych.}
              Graf wiedzy budujemy na podstawie ręcznie przygotowanych plików CSV; brak automatycznej
              synchronizacji z corocznymi zmianami regulaminów i komunikatów rektorskich.
              System nie informuje użytkownika o dacie ważności pobranych danych.

        \item \textbf{Weryfikacja faktów oparta na tym samym modelu.}
              Audytor faktów korzysta z tego samego LLM co generator, co nie gwarantuje niezależnej kontroli.
              Weryfikacja może być zbyt restrykcyjna (odrzucenie poprawnego szkicu) lub niewystarczająca ---
              np.\ przy zapytaniu o przeniesienie na Inżynierię Materiałową oba systemy hybrydowe podały
              poprawny stosunek $176/60 \approx 2{,}93$, lecz błędnie sformułowały go jako „176 miejsc na 60 kandydatów''.

        \item \textbf{Brak warstwy użytkowej.}
              Obecna implementacja to narzędzie deweloperskie (CLI, skrypty testowe), bez interfejsu
              przeznaczonego dla kandydata lub studenta AGH. System nie jest gotowy do wdrożenia produkcyjnego.

        \item \textbf{Ograniczona skala testów.}
              Zestaw 18 scenariuszy pozwala na orientacyjne porównanie architektur, lecz nie stanowi
              pełnej walidacji statystycznej. Brakuje automatycznego scoringu względem referencyjnych odpowiedzi.
    \end{itemize}
\end{tcolorbox}

\vspace{1em}

% ---------------------------------------------------------------------------
\subsection*{Kierunki dalszego rozwoju}

\begin{tcolorbox}[colback=green!5!white,colframe=green!50!black,title=Proponowane kolejne kroki]
    \begin{enumerate}[leftmargin=*, nosep]
        \item \textbf{Rozbudowa i utrzymanie bazy wiedzy.}
              Automatyzacja ingestu dokumentów PDF, wersjonowanie źródeł według roku akademickiego
              oraz uzupełnienie grafu o brakujące encje (stypendia, komunikaty rektorskie, progi bieżącej rekrutacji).

        \item \textbf{Ulepszenie schematu grafu.}
              Jawne modelowanie relacji logicznych (alternatywa, koniunkcja) w węzłach przedmiotów maturalnych
              i warunków rekrutacyjnych, aby uniknąć spłaszczania reguł tekstowych do list encji.

        \item \textbf{Ważenie i fuzja wyników pobierania.}
              Wprowadzenie mechanizmu rerankingu zamiast prostego łączenia kontekstu grafowego i wektorowego,
              z możliwością dynamicznego priorytetyzowania źródła w zależności od rozpoznanej intencji.

        \item \textbf{Obiektywna metryka ewaluacji.}
              Opracowanie zestawu referencyjnych odpowiedzi (gold standard) i automatycznego scoringu
              (np.\ pokrycie faktów, poprawność liczbowa), uzupełniającego obecną ocenę jakościową.

        \item \textbf{Optymalizacja czasu reakcji.}
              Buforowanie wyników analizy intencji, równoległe wywołania niezależnych agentów
              oraz zastosowanie lżejszego modelu do zadań klasyfikacyjnych (strażnik, filtr).

        \item \textbf{Warstwa prezentacji.}
              Prosty interfejs webowy lub integracja z istniejącymi kanałami informacyjnymi uczelni,
              z podaniem źródeł i daty aktualności danych.

        \item \textbf{Wzmocnienie zabezpieczeń.}
              Rozszerzenie zestawu testów obejścia (fałszywe paragrafy regulaminu, wymuszanie halucynacji)
              oraz uzupełnienie klasyfikacji LLM o reguły deterministyczne dla znanych wzorców ataku.
    \end{enumerate}
\end{tcolorbox}

\vspace{1em}

\noindent
Podsumowując, projekt zrealizował założony cel porównawczy trzech architektur RAG w domenie uczelnianej.
Największą wartość edukacyjną --- zarówno dla nas jako zespołu studenckiego, jak i dla przyszłych prac
nad systemem --- stanowi praktyczne zrozumienie kompromisów między prostotą implementacji, precyzją
danych strukturalnych a wieloetapową kontrolą jakości odpowiedzi. Żadne pojedyncze podejście nie rozwiązuje
wszystkich zidentyfikowanych problemów; sensownym kierunkiem rozwoju wydaje się dalsze pogłębianie warstwy
danych i ewaluacji, przy stopniowym włączaniu mechanizmów agentowych tam, gdzie wiarygodność odpowiedzi
ma pierwszeństwo przed minimalnym czasem reakcji.

\vspace{1em}

\newpage

% Wymaga pakietów z main.tex: tcolorbox, enumitem

\section*{Podsumowanie i refleksja projektowa}

Celem projektu było zaprojektowanie i porównanie trzech wariantów asystenta uczelnianego AGH
opartego na generowaniu wspomaganym pobieraniem informacji (RAG): wyszukiwania wyłącznie wektorowego,
hybrydowego GraphRAG (graf Neo4j + wektory Chroma) oraz wieloetapowego pipeline'u agentowego (LangGraph).
Ocenę przeprowadziliśmy na zestawie 18 scenariuszy --- 13 pytań merytorycznych dotyczących rekrutacji
i studiów oraz 5 prób obejścia zabezpieczeń (manipulacja promptem, fałszywe źródła, żądania nielegalne).
Dla każdego pytania zarejestrowaliśmy odpowiedzi trzech systemów, czas reakcji oraz --- w wariancie agentowym ---
pełny ślad etapów przetwarzania (trace JSON).

\vspace{1em}

% ---------------------------------------------------------------------------
\subsection*{Wnioski z pracy zespołowej}

\begin{tcolorbox}[colback=blue!5!white,colframe=blue!75!black,title=Doświadczenia edukacyjne zespołu]
    \begin{itemize}[leftmargin=*, nosep]
        \item \textbf{Projektowanie bazy wiedzy hybrydowej.}
              Zrozumieliśmy, że połączenie danych strukturalnych (tabele CSV załadowane do Neo4j)
              z tekstem regulaminów (fragmenty PDF w Chroma) wymaga świadomego modelowania ontologii ---
              relacji między kierunkiem, progiem punktowym, przedmiotem maturalnym czy terminem rekrutacji.
              Jakość odpowiedzi zależy w równym stopniu od schematu grafu i jakości wektoryzacji dokumentów.

        \item \textbf{Brak uniwersalnie najlepszej architektury.}
              Wyniki benchmarku potwierdziły, że żaden wariant nie dominuje we wszystkich kategoriach zapytań.
              Wyszukiwanie wektorowe dobrze odtwarza fragmenty tekstu z zachowaną logiką słowną
              (np.\ operator „albo'' w tabeli przedmiotów), lecz zawodzi przy filtrowaniu numerycznym
              (progi poniżej 700 punktów) i zapytaniach wielowarunkowych (terminy konkretnego cyklu rekrutacji).
              GraphRAG precyzyjnie obsługuje zapytania relacyjne i arytmetyczne, natomiast płaski zapis encji
              w grafie może zniekształcić reguły logiczne źródłowe.

        \item \textbf{Orkiestracja wieloetapowa jako warstwa kontroli.}
              Budowa grafu stanów w LangGraph nauczyła nas rozdzielać odpowiedzialności między wyspecjalizowane
              moduły: strażnika bezpieczeństwa, analizator intencji, filtr kontekstu, generator odpowiedzi
              i audytora faktów. Każdy etap podejmuje decyzję warunkową (odmowa, doprecyzowanie, akceptacja),
              co odróżnia system agentowy od liniowego schematu \textit{pobierz\,--\,wygeneruj}.

        \item \textbf{Konieczność systemowej ewaluacji.}
              Opracowanie narzędzia porównawczego (\texttt{run\_compare.py}) i rejestrowanie śladów JSON
              uświadomiło nam, że ocena asystenta nie może ograniczać się do poprawności merytorycznej końcowej odpowiedzi.
              Istotne są także: etap, na którym system zawiódł, czas reakcji, zachowanie przy braku danych
              oraz sposób reagowania na próby manipulacji.

        \item \textbf{Praca zespołowa w projekcie inżynierskim.}
              Podzieliliśmy zakres na komplementarne obszary: przygotowanie i ingest danych, moduły pobierania
              informacji, zamrożone pipeline'y porównawcze, warstwę agentową oraz dokumentację z benchmarkiem.
              Integracja komponentów wymagała uzgodnienia wspólnych interfejsów stanu i formatów kontekstu,
              co stanowiło istotny element nauki inżynierii oprogramowania na rzeczywistym przypadku użycia.
    \end{itemize}
\end{tcolorbox}

\vspace{1em}

% ---------------------------------------------------------------------------
\subsection*{Zidentyfikowane słabości}

\begin{tcolorbox}[colback=red!5!white,colframe=red!50!black,title=Ograniczenia obecnej implementacji]
    \begin{itemize}[leftmargin=*, nosep]
        \item \textbf{Niejednorodna jakość odpowiedzi.}
              W zależności od kategorii pytania każdy wariant wykazuje inne słabości.
              Przykład: przy pytaniu o punkty za osiągnięcia sportowe wyszukiwanie wektorowe zwróciło fragment
              o punktach ECTS zamiast tabeli rekrutacyjnej; GraphRAG pomylił tabele (odznaki PTTK zamiast klas sportowych);
              system agentowy odrzucił błędny szkic w weryfikacji faktów, lecz użytkownik nie otrzymał pełnej,
              poprawnej odpowiedzi z listą dyscyplin.

        \item \textbf{Ograniczenia modelu grafowego.}
              Przy przedmiotach maturalnych na kierunku Automatyka i Robotyka graf zwrócił płaską listę encji
              bez operatora logicznego „albo'', co doprowadziło GraphRAG i RAG agentowy do błędnej interpretacji
              (wymaganie wszystkich przedmiotów naraz). Przy pytaniu o stypendium socjalne graf nie zwrócił
              żadnych wyników, a system oparł się wyłącznie na słabych fragmentach wektorowych.

        \item \textbf{Wysoki koszt obliczeniowy wariantu agentowego.}
              Pipeline agentowy wymaga od czterech do sześciu wywołań modelu językowego na jedno pytanie,
              co przekłada się na wielokrotnie wyższą latencję w porównaniu z wariantami liniowymi.
              W scenariuszach interaktywnych może to obniżać użyteczność systemu.

        \item \textbf{Zależność od jakości i aktualności danych wejściowych.}
              Graf wiedzy budujemy na podstawie ręcznie przygotowanych plików CSV; brak automatycznej
              synchronizacji z corocznymi zmianami regulaminów i komunikatów rektorskich.
              System nie informuje użytkownika o dacie ważności pobranych danych.

        \item \textbf{Weryfikacja faktów oparta na tym samym modelu.}
              Audytor faktów korzysta z tego samego LLM co generator, co nie gwarantuje niezależnej kontroli.
              Weryfikacja może być zbyt restrykcyjna (odrzucenie poprawnego szkicu) lub niewystarczająca ---
              np.\ przy zapytaniu o przeniesienie na Inżynierię Materiałową oba systemy hybrydowe podały
              poprawny stosunek $176/60 \approx 2{,}93$, lecz błędnie sformułowały go jako „176 miejsc na 60 kandydatów''.

        \item \textbf{Brak warstwy użytkowej.}
              Obecna implementacja to narzędzie deweloperskie (CLI, skrypty testowe), bez interfejsu
              przeznaczonego dla kandydata lub studenta AGH. System nie jest gotowy do wdrożenia produkcyjnego.

        \item \textbf{Ograniczona skala testów.}
              Zestaw 18 scenariuszy pozwala na orientacyjne porównanie architektur, lecz nie stanowi
              pełnej walidacji statystycznej. Brakuje automatycznego scoringu względem referencyjnych odpowiedzi.
    \end{itemize}
\end{tcolorbox}

\vspace{1em}

% ---------------------------------------------------------------------------
\subsection*{Kierunki dalszego rozwoju}

\begin{tcolorbox}[colback=green!5!white,colframe=green!50!black,title=Proponowane kolejne kroki]
    \begin{enumerate}[leftmargin=*, nosep]
        \item \textbf{Rozbudowa i utrzymanie bazy wiedzy.}
              Automatyzacja ingestu dokumentów PDF, wersjonowanie źródeł według roku akademickiego
              oraz uzupełnienie grafu o brakujące encje (stypendia, komunikaty rektorskie, progi bieżącej rekrutacji).

        \item \textbf{Ulepszenie schematu grafu.}
              Jawne modelowanie relacji logicznych (alternatywa, koniunkcja) w węzłach przedmiotów maturalnych
              i warunków rekrutacyjnych, aby uniknąć spłaszczania reguł tekstowych do list encji.

        \item \textbf{Ważenie i fuzja wyników pobierania.}
              Wprowadzenie mechanizmu rerankingu zamiast prostego łączenia kontekstu grafowego i wektorowego,
              z możliwością dynamicznego priorytetyzowania źródła w zależności od rozpoznanej intencji.

        \item \textbf{Obiektywna metryka ewaluacji.}
              Opracowanie zestawu referencyjnych odpowiedzi (gold standard) i automatycznego scoringu
              (np.\ pokrycie faktów, poprawność liczbowa), uzupełniającego obecną ocenę jakościową.

        \item \textbf{Optymalizacja czasu reakcji.}
              Buforowanie wyników analizy intencji, równoległe wywołania niezależnych agentów
              oraz zastosowanie lżejszego modelu do zadań klasyfikacyjnych (strażnik, filtr).

        \item \textbf{Warstwa prezentacji.}
              Prosty interfejs webowy lub integracja z istniejącymi kanałami informacyjnymi uczelni,
              z podaniem źródeł i daty aktualności danych.

        \item \textbf{Wzmocnienie zabezpieczeń.}
              Rozszerzenie zestawu testów obejścia (fałszywe paragrafy regulaminu, wymuszanie halucynacji)
              oraz uzupełnienie klasyfikacji LLM o reguły deterministyczne dla znanych wzorców ataku.
    \end{enumerate}
\end{tcolorbox}

\vspace{1em}

\noindent
Podsumowując, projekt zrealizował założony cel porównawczy trzech architektur RAG w domenie uczelnianej.
Największą wartość edukacyjną --- zarówno dla nas jako zespołu studenckiego, jak i dla przyszłych prac
nad systemem --- stanowi praktyczne zrozumienie kompromisów między prostotą implementacji, precyzją
danych strukturalnych a wieloetapową kontrolą jakości odpowiedzi. Żadne pojedyncze podejście nie rozwiązuje
wszystkich zidentyfikowanych problemów; sensownym kierunkiem rozwoju wydaje się dalsze pogłębianie warstwy
danych i ewaluacji, przy stopniowym włączaniu mechanizmów agentowych tam, gdzie wiarygodność odpowiedzi
ma pierwszeństwo przed minimalnym czasem reakcji.

\vspace{1em}

\newpage

% Wymaga pakietów z main.tex: tcolorbox, enumitem

\section*{Podsumowanie i refleksja projektowa}

Celem projektu było zaprojektowanie i porównanie trzech wariantów asystenta uczelnianego AGH
opartego na generowaniu wspomaganym pobieraniem informacji (RAG): wyszukiwania wyłącznie wektorowego,
hybrydowego GraphRAG (graf Neo4j + wektory Chroma) oraz wieloetapowego pipeline'u agentowego (LangGraph).
Ocenę przeprowadziliśmy na zestawie 18 scenariuszy --- 13 pytań merytorycznych dotyczących rekrutacji
i studiów oraz 5 prób obejścia zabezpieczeń (manipulacja promptem, fałszywe źródła, żądania nielegalne).
Dla każdego pytania zarejestrowaliśmy odpowiedzi trzech systemów, czas reakcji oraz --- w wariancie agentowym ---
pełny ślad etapów przetwarzania (trace JSON).

\vspace{1em}

% ---------------------------------------------------------------------------
\subsection*{Wnioski z pracy zespołowej}

\begin{tcolorbox}[colback=blue!5!white,colframe=blue!75!black,title=Doświadczenia edukacyjne zespołu]
    \begin{itemize}[leftmargin=*, nosep]
        \item \textbf{Projektowanie bazy wiedzy hybrydowej.}
              Zrozumieliśmy, że połączenie danych strukturalnych (tabele CSV załadowane do Neo4j)
              z tekstem regulaminów (fragmenty PDF w Chroma) wymaga świadomego modelowania ontologii ---
              relacji między kierunkiem, progiem punktowym, przedmiotem maturalnym czy terminem rekrutacji.
              Jakość odpowiedzi zależy w równym stopniu od schematu grafu i jakości wektoryzacji dokumentów.

        \item \textbf{Brak uniwersalnie najlepszej architektury.}
              Wyniki benchmarku potwierdziły, że żaden wariant nie dominuje we wszystkich kategoriach zapytań.
              Wyszukiwanie wektorowe dobrze odtwarza fragmenty tekstu z zachowaną logiką słowną
              (np.\ operator „albo'' w tabeli przedmiotów), lecz zawodzi przy filtrowaniu numerycznym
              (progi poniżej 700 punktów) i zapytaniach wielowarunkowych (terminy konkretnego cyklu rekrutacji).
              GraphRAG precyzyjnie obsługuje zapytania relacyjne i arytmetyczne, natomiast płaski zapis encji
              w grafie może zniekształcić reguły logiczne źródłowe.

        \item \textbf{Orkiestracja wieloetapowa jako warstwa kontroli.}
              Budowa grafu stanów w LangGraph nauczyła nas rozdzielać odpowiedzialności między wyspecjalizowane
              moduły: strażnika bezpieczeństwa, analizator intencji, filtr kontekstu, generator odpowiedzi
              i audytora faktów. Każdy etap podejmuje decyzję warunkową (odmowa, doprecyzowanie, akceptacja),
              co odróżnia system agentowy od liniowego schematu \textit{pobierz\,--\,wygeneruj}.

        \item \textbf{Konieczność systemowej ewaluacji.}
              Opracowanie narzędzia porównawczego (\texttt{run\_compare.py}) i rejestrowanie śladów JSON
              uświadomiło nam, że ocena asystenta nie może ograniczać się do poprawności merytorycznej końcowej odpowiedzi.
              Istotne są także: etap, na którym system zawiódł, czas reakcji, zachowanie przy braku danych
              oraz sposób reagowania na próby manipulacji.

        \item \textbf{Praca zespołowa w projekcie inżynierskim.}
              Podzieliliśmy zakres na komplementarne obszary: przygotowanie i ingest danych, moduły pobierania
              informacji, zamrożone pipeline'y porównawcze, warstwę agentową oraz dokumentację z benchmarkiem.
              Integracja komponentów wymagała uzgodnienia wspólnych interfejsów stanu i formatów kontekstu,
              co stanowiło istotny element nauki inżynierii oprogramowania na rzeczywistym przypadku użycia.
    \end{itemize}
\end{tcolorbox}

\vspace{1em}

% ---------------------------------------------------------------------------
\subsection*{Zidentyfikowane słabości}

\begin{tcolorbox}[colback=red!5!white,colframe=red!50!black,title=Ograniczenia obecnej implementacji]
    \begin{itemize}[leftmargin=*, nosep]
        \item \textbf{Niejednorodna jakość odpowiedzi.}
              W zależności od kategorii pytania każdy wariant wykazuje inne słabości.
              Przykład: przy pytaniu o punkty za osiągnięcia sportowe wyszukiwanie wektorowe zwróciło fragment
              o punktach ECTS zamiast tabeli rekrutacyjnej; GraphRAG pomylił tabele (odznaki PTTK zamiast klas sportowych);
              system agentowy odrzucił błędny szkic w weryfikacji faktów, lecz użytkownik nie otrzymał pełnej,
              poprawnej odpowiedzi z listą dyscyplin.

        \item \textbf{Ograniczenia modelu grafowego.}
              Przy przedmiotach maturalnych na kierunku Automatyka i Robotyka graf zwrócił płaską listę encji
              bez operatora logicznego „albo'', co doprowadziło GraphRAG i RAG agentowy do błędnej interpretacji
              (wymaganie wszystkich przedmiotów naraz). Przy pytaniu o stypendium socjalne graf nie zwrócił
              żadnych wyników, a system oparł się wyłącznie na słabych fragmentach wektorowych.

        \item \textbf{Wysoki koszt obliczeniowy wariantu agentowego.}
              Pipeline agentowy wymaga od czterech do sześciu wywołań modelu językowego na jedno pytanie,
              co przekłada się na wielokrotnie wyższą latencję w porównaniu z wariantami liniowymi.
              W scenariuszach interaktywnych może to obniżać użyteczność systemu.

        \item \textbf{Zależność od jakości i aktualności danych wejściowych.}
              Graf wiedzy budujemy na podstawie ręcznie przygotowanych plików CSV; brak automatycznej
              synchronizacji z corocznymi zmianami regulaminów i komunikatów rektorskich.
              System nie informuje użytkownika o dacie ważności pobranych danych.

        \item \textbf{Weryfikacja faktów oparta na tym samym modelu.}
              Audytor faktów korzysta z tego samego LLM co generator, co nie gwarantuje niezależnej kontroli.
              Weryfikacja może być zbyt restrykcyjna (odrzucenie poprawnego szkicu) lub niewystarczająca ---
              np.\ przy zapytaniu o przeniesienie na Inżynierię Materiałową oba systemy hybrydowe podały
              poprawny stosunek $176/60 \approx 2{,}93$, lecz błędnie sformułowały go jako „176 miejsc na 60 kandydatów''.

        \item \textbf{Brak warstwy użytkowej.}
              Obecna implementacja to narzędzie deweloperskie (CLI, skrypty testowe), bez interfejsu
              przeznaczonego dla kandydata lub studenta AGH. System nie jest gotowy do wdrożenia produkcyjnego.

        \item \textbf{Ograniczona skala testów.}
              Zestaw 18 scenariuszy pozwala na orientacyjne porównanie architektur, lecz nie stanowi
              pełnej walidacji statystycznej. Brakuje automatycznego scoringu względem referencyjnych odpowiedzi.
    \end{itemize}
\end{tcolorbox}

\vspace{1em}

% ---------------------------------------------------------------------------
\subsection*{Kierunki dalszego rozwoju}

\begin{tcolorbox}[colback=green!5!white,colframe=green!50!black,title=Proponowane kolejne kroki]
    \begin{enumerate}[leftmargin=*, nosep]
        \item \textbf{Rozbudowa i utrzymanie bazy wiedzy.}
              Automatyzacja ingestu dokumentów PDF, wersjonowanie źródeł według roku akademickiego
              oraz uzupełnienie grafu o brakujące encje (stypendia, komunikaty rektorskie, progi bieżącej rekrutacji).

        \item \textbf{Ulepszenie schematu grafu.}
              Jawne modelowanie relacji logicznych (alternatywa, koniunkcja) w węzłach przedmiotów maturalnych
              i warunków rekrutacyjnych, aby uniknąć spłaszczania reguł tekstowych do list encji.

        \item \textbf{Ważenie i fuzja wyników pobierania.}
              Wprowadzenie mechanizmu rerankingu zamiast prostego łączenia kontekstu grafowego i wektorowego,
              z możliwością dynamicznego priorytetyzowania źródła w zależności od rozpoznanej intencji.

        \item \textbf{Obiektywna metryka ewaluacji.}
              Opracowanie zestawu referencyjnych odpowiedzi (gold standard) i automatycznego scoringu
              (np.\ pokrycie faktów, poprawność liczbowa), uzupełniającego obecną ocenę jakościową.

        \item \textbf{Optymalizacja czasu reakcji.}
              Buforowanie wyników analizy intencji, równoległe wywołania niezależnych agentów
              oraz zastosowanie lżejszego modelu do zadań klasyfikacyjnych (strażnik, filtr).

        \item \textbf{Warstwa prezentacji.}
              Prosty interfejs webowy lub integracja z istniejącymi kanałami informacyjnymi uczelni,
              z podaniem źródeł i daty aktualności danych.

        \item \textbf{Wzmocnienie zabezpieczeń.}
              Rozszerzenie zestawu testów obejścia (fałszywe paragrafy regulaminu, wymuszanie halucynacji)
              oraz uzupełnienie klasyfikacji LLM o reguły deterministyczne dla znanych wzorców ataku.
    \end{enumerate}
\end{tcolorbox}

\vspace{1em}

\noindent
Podsumowując, projekt zrealizował założony cel porównawczy trzech architektur RAG w domenie uczelnianej.
Największą wartość edukacyjną --- zarówno dla nas jako zespołu studenckiego, jak i dla przyszłych prac
nad systemem --- stanowi praktyczne zrozumienie kompromisów między prostotą implementacji, precyzją
danych strukturalnych a wieloetapową kontrolą jakości odpowiedzi. Żadne pojedyncze podejście nie rozwiązuje
wszystkich zidentyfikowanych problemów; sensownym kierunkiem rozwoju wydaje się dalsze pogłębianie warstwy
danych i ewaluacji, przy stopniowym włączaniu mechanizmów agentowych tam, gdzie wiarygodność odpowiedzi
ma pierwszeństwo przed minimalnym czasem reakcji.

\vspace{1em}

\newpage

% Wymaga pakietów z main.tex: tcolorbox, enumitem

\section*{Podsumowanie i refleksja projektowa}

Celem projektu było zaprojektowanie i porównanie trzech wariantów asystenta uczelnianego AGH
opartego na generowaniu wspomaganym pobieraniem informacji (RAG): wyszukiwania wyłącznie wektorowego,
hybrydowego GraphRAG (graf Neo4j + wektory Chroma) oraz wieloetapowego pipeline'u agentowego (LangGraph).
Ocenę przeprowadziliśmy na zestawie 18 scenariuszy --- 13 pytań merytorycznych dotyczących rekrutacji
i studiów oraz 5 prób obejścia zabezpieczeń (manipulacja promptem, fałszywe źródła, żądania nielegalne).
Dla każdego pytania zarejestrowaliśmy odpowiedzi trzech systemów, czas reakcji oraz --- w wariancie agentowym ---
pełny ślad etapów przetwarzania (trace JSON).

\vspace{1em}

% ---------------------------------------------------------------------------
\subsection*{Wnioski z pracy zespołowej}

\begin{tcolorbox}[colback=blue!5!white,colframe=blue!75!black,title=Doświadczenia edukacyjne zespołu]
    \begin{itemize}[leftmargin=*, nosep]
        \item \textbf{Projektowanie bazy wiedzy hybrydowej.}
              Zrozumieliśmy, że połączenie danych strukturalnych (tabele CSV załadowane do Neo4j)
              z tekstem regulaminów (fragmenty PDF w Chroma) wymaga świadomego modelowania ontologii ---
              relacji między kierunkiem, progiem punktowym, przedmiotem maturalnym czy terminem rekrutacji.
              Jakość odpowiedzi zależy w równym stopniu od schematu grafu i jakości wektoryzacji dokumentów.

        \item \textbf{Brak uniwersalnie najlepszej architektury.}
              Wyniki benchmarku potwierdziły, że żaden wariant nie dominuje we wszystkich kategoriach zapytań.
              Wyszukiwanie wektorowe dobrze odtwarza fragmenty tekstu z zachowaną logiką słowną
              (np.\ operator „albo'' w tabeli przedmiotów), lecz zawodzi przy filtrowaniu numerycznym
              (progi poniżej 700 punktów) i zapytaniach wielowarunkowych (terminy konkretnego cyklu rekrutacji).
              GraphRAG precyzyjnie obsługuje zapytania relacyjne i arytmetyczne, natomiast płaski zapis encji
              w grafie może zniekształcić reguły logiczne źródłowe.

        \item \textbf{Orkiestracja wieloetapowa jako warstwa kontroli.}
              Budowa grafu stanów w LangGraph nauczyła nas rozdzielać odpowiedzialności między wyspecjalizowane
              moduły: strażnika bezpieczeństwa, analizator intencji, filtr kontekstu, generator odpowiedzi
              i audytora faktów. Każdy etap podejmuje decyzję warunkową (odmowa, doprecyzowanie, akceptacja),
              co odróżnia system agentowy od liniowego schematu \textit{pobierz\,--\,wygeneruj}.

        \item \textbf{Konieczność systemowej ewaluacji.}
              Opracowanie narzędzia porównawczego (\texttt{run\_compare.py}) i rejestrowanie śladów JSON
              uświadomiło nam, że ocena asystenta nie może ograniczać się do poprawności merytorycznej końcowej odpowiedzi.
              Istotne są także: etap, na którym system zawiódł, czas reakcji, zachowanie przy braku danych
              oraz sposób reagowania na próby manipulacji.

        \item \textbf{Praca zespołowa w projekcie inżynierskim.}
              Podzieliliśmy zakres na komplementarne obszary: przygotowanie i ingest danych, moduły pobierania
              informacji, zamrożone pipeline'y porównawcze, warstwę agentową oraz dokumentację z benchmarkiem.
              Integracja komponentów wymagała uzgodnienia wspólnych interfejsów stanu i formatów kontekstu,
              co stanowiło istotny element nauki inżynierii oprogramowania na rzeczywistym przypadku użycia.
    \end{itemize}
\end{tcolorbox}

\vspace{1em}

% ---------------------------------------------------------------------------
\subsection*{Zidentyfikowane słabości}

\begin{tcolorbox}[colback=red!5!white,colframe=red!50!black,title=Ograniczenia obecnej implementacji]
    \begin{itemize}[leftmargin=*, nosep]
        \item \textbf{Niejednorodna jakość odpowiedzi.}
              W zależności od kategorii pytania każdy wariant wykazuje inne słabości.
              Przykład: przy pytaniu o punkty za osiągnięcia sportowe wyszukiwanie wektorowe zwróciło fragment
              o punktach ECTS zamiast tabeli rekrutacyjnej; GraphRAG pomylił tabele (odznaki PTTK zamiast klas sportowych);
              system agentowy odrzucił błędny szkic w weryfikacji faktów, lecz użytkownik nie otrzymał pełnej,
              poprawnej odpowiedzi z listą dyscyplin.

        \item \textbf{Ograniczenia modelu grafowego.}
              Przy przedmiotach maturalnych na kierunku Automatyka i Robotyka graf zwrócił płaską listę encji
              bez operatora logicznego „albo'', co doprowadziło GraphRAG i RAG agentowy do błędnej interpretacji
              (wymaganie wszystkich przedmiotów naraz). Przy pytaniu o stypendium socjalne graf nie zwrócił
              żadnych wyników, a system oparł się wyłącznie na słabych fragmentach wektorowych.

        \item \textbf{Wysoki koszt obliczeniowy wariantu agentowego.}
              Pipeline agentowy wymaga od czterech do sześciu wywołań modelu językowego na jedno pytanie,
              co przekłada się na wielokrotnie wyższą latencję w porównaniu z wariantami liniowymi.
              W scenariuszach interaktywnych może to obniżać użyteczność systemu.

        \item \textbf{Zależność od jakości i aktualności danych wejściowych.}
              Graf wiedzy budujemy na podstawie ręcznie przygotowanych plików CSV; brak automatycznej
              synchronizacji z corocznymi zmianami regulaminów i komunikatów rektorskich.
              System nie informuje użytkownika o dacie ważności pobranych danych.

        \item \textbf{Weryfikacja faktów oparta na tym samym modelu.}
              Audytor faktów korzysta z tego samego LLM co generator, co nie gwarantuje niezależnej kontroli.
              Weryfikacja może być zbyt restrykcyjna (odrzucenie poprawnego szkicu) lub niewystarczająca ---
              np.\ przy zapytaniu o przeniesienie na Inżynierię Materiałową oba systemy hybrydowe podały
              poprawny stosunek $176/60 \approx 2{,}93$, lecz błędnie sformułowały go jako „176 miejsc na 60 kandydatów''.

        \item \textbf{Brak warstwy użytkowej.}
              Obecna implementacja to narzędzie deweloperskie (CLI, skrypty testowe), bez interfejsu
              przeznaczonego dla kandydata lub studenta AGH. System nie jest gotowy do wdrożenia produkcyjnego.

        \item \textbf{Ograniczona skala testów.}
              Zestaw 18 scenariuszy pozwala na orientacyjne porównanie architektur, lecz nie stanowi
              pełnej walidacji statystycznej. Brakuje automatycznego scoringu względem referencyjnych odpowiedzi.
    \end{itemize}
\end{tcolorbox}

\vspace{1em}

% ---------------------------------------------------------------------------
\subsection*{Kierunki dalszego rozwoju}

\begin{tcolorbox}[colback=green!5!white,colframe=green!50!black,title=Proponowane kolejne kroki]
    \begin{enumerate}[leftmargin=*, nosep]
        \item \textbf{Rozbudowa i utrzymanie bazy wiedzy.}
              Automatyzacja ingestu dokumentów PDF, wersjonowanie źródeł według roku akademickiego
              oraz uzupełnienie grafu o brakujące encje (stypendia, komunikaty rektorskie, progi bieżącej rekrutacji).

        \item \textbf{Ulepszenie schematu grafu.}
              Jawne modelowanie relacji logicznych (alternatywa, koniunkcja) w węzłach przedmiotów maturalnych
              i warunków rekrutacyjnych, aby uniknąć spłaszczania reguł tekstowych do list encji.

        \item \textbf{Ważenie i fuzja wyników pobierania.}
              Wprowadzenie mechanizmu rerankingu zamiast prostego łączenia kontekstu grafowego i wektorowego,
              z możliwością dynamicznego priorytetyzowania źródła w zależności od rozpoznanej intencji.

        \item \textbf{Obiektywna metryka ewaluacji.}
              Opracowanie zestawu referencyjnych odpowiedzi (gold standard) i automatycznego scoringu
              (np.\ pokrycie faktów, poprawność liczbowa), uzupełniającego obecną ocenę jakościową.

        \item \textbf{Optymalizacja czasu reakcji.}
              Buforowanie wyników analizy intencji, równoległe wywołania niezależnych agentów
              oraz zastosowanie lżejszego modelu do zadań klasyfikacyjnych (strażnik, filtr).

        \item \textbf{Warstwa prezentacji.}
              Prosty interfejs webowy lub integracja z istniejącymi kanałami informacyjnymi uczelni,
              z podaniem źródeł i daty aktualności danych.

        \item \textbf{Wzmocnienie zabezpieczeń.}
              Rozszerzenie zestawu testów obejścia (fałszywe paragrafy regulaminu, wymuszanie halucynacji)
              oraz uzupełnienie klasyfikacji LLM o reguły deterministyczne dla znanych wzorców ataku.
    \end{enumerate}
\end{tcolorbox}

\vspace{1em}

\noindent
Podsumowując, projekt zrealizował założony cel porównawczy trzech architektur RAG w domenie uczelnianej.
Największą wartość edukacyjną --- zarówno dla nas jako zespołu studenckiego, jak i dla przyszłych prac
nad systemem --- stanowi praktyczne zrozumienie kompromisów między prostotą implementacji, precyzją
danych strukturalnych a wieloetapową kontrolą jakości odpowiedzi. Żadne pojedyncze podejście nie rozwiązuje
wszystkich zidentyfikowanych problemów; sensownym kierunkiem rozwoju wydaje się dalsze pogłębianie warstwy
danych i ewaluacji, przy stopniowym włączaniu mechanizmów agentowych tam, gdzie wiarygodność odpowiedzi
ma pierwszeństwo przed minimalnym czasem reakcji.

\vspace{1em}

\newpage
